\documentclass[11pt,a4paper]{article}
\usepackage{amsmath}
\usepackage{amssymb}
\usepackage{geometry}
\usepackage{graphicx}
\usepackage{hyperref}
\usepackage{booktabs}
\usepackage{bm}
\usepackage{siunitx}

\geometry{margin=1in}

\title{Sensor Systems and Vehicle Dynamics:\\
Mathematical Foundations and Theory}
\author{Mario Tilocca}
\date{\today}

\begin{document}

\maketitle

\begin{abstract}
This document provides the complete mathematical foundation for the sensor simulation and vehicle dynamics model used in autonomous mining vehicle development. We derive the bicycle model kinematics, Ackermann steering geometry, longitudinal dynamics with energy flow, and comprehensive sensor noise models including IMU bias dynamics, GNSS random walk, and coordinate transformations between vehicle body frame, NED frame, and WGS84 geodetic coordinates.
\end{abstract}

\tableofcontents
\clearpage

\section{Coordinate Systems and Transformations}

\subsection{Reference Frames}

\subsubsection{Global Inertial Frame (World)}

Fixed reference frame with origin at simulation start:
\begin{itemize}
    \item $X$ -- East (or initial vehicle forward direction)
    \item $Y$ -- North (or initial vehicle left direction)
    \item $Z$ -- Up (vertical, right-hand rule)
\end{itemize}

Vehicle state in global frame: $(x, y, \psi)$ where $\psi$ is heading angle (yaw) measured counterclockwise from $X$-axis.

\subsubsection{Vehicle Body Frame}

Moving frame attached to vehicle with origin at rear axle center:
\begin{itemize}
    \item $x_b$ -- Forward (longitudinal)
    \item $y_b$ -- Left (lateral)
    \item $z_b$ -- Up (vertical)
\end{itemize}

Vehicle velocity in body frame: $v$ (longitudinal speed, m/s).

\subsubsection{North-East-Down (NED) Frame}

Navigation frame used by GNSS:
\begin{itemize}
    \item $N$ -- North
    \item $E$ -- East
    \item $D$ -- Down (toward Earth center)
\end{itemize}

\subsection{Coordinate Transformations}

\subsubsection{Body to Global}

Position rate of change:
\begin{align}
\dot{x} &= v \cos(\psi) \label{eq:xdot}\\
\dot{y} &= v \sin(\psi) \label{eq:ydot}
\end{align}

Rotation matrix from body to global:
\begin{equation}
\mathbf{R}_b^g = \begin{bmatrix}
\cos\psi & -\sin\psi & 0 \\
\sin\psi & \cos\psi & 0 \\
0 & 0 & 1
\end{bmatrix}
\end{equation}

\subsubsection{Body to NED}

Velocity transformation:
\begin{align}
v_N &= v \sin(\psi) \\
v_E &= v \cos(\psi) \\
v_D &= 0 \quad \text{(ground vehicle)}
\end{align}

\subsubsection{Local XY to WGS84}

Flat Earth approximation for short distances ($<10$ km):
\begin{align}
\text{meters per degree latitude:} \quad m_{\phi} &= 111{,}320 \text{ m} \\
\text{meters per degree longitude:} \quad m_{\lambda} &= 111{,}320 \cdot \cos(\phi_0) \text{ m}
\end{align}

Position conversion:
\begin{align}
\phi &= \phi_0 + \frac{y}{m_{\phi}} \\
\lambda &= \lambda_0 + \frac{x}{m_{\lambda}}
\end{align}

where $(\phi_0, \lambda_0)$ is the origin in WGS84 coordinates (latitude, longitude in degrees).

\section{Bicycle Model Kinematics}

\subsection{Model Assumptions}

\begin{enumerate}
    \item \textbf{No slip:} Pure rolling condition (wheels do not slide)
    \item \textbf{Planar motion:} Vehicle remains on flat ground (no pitch/roll)
    \item \textbf{Point mass:} Neglect rotational inertia about vertical axis
    \item \textbf{Rear-axle reference:} Kinematic center at rear axle midpoint
\end{enumerate}

\subsection{Kinematic Equations}

The bicycle model reduces a four-wheeled vehicle to two wheels separated by wheelbase $L$:

\begin{equation}
\begin{aligned}
\dot{x} &= v \cos(\psi) \\
\dot{y} &= v \sin(\psi) \\
\dot{\psi} &= \frac{v}{L} \tan(\delta)
\end{aligned}
\label{eq:bicycle_model}
\end{equation}

where:
\begin{itemize}
    \item $(x, y)$ -- position of rear axle center in global frame (m)
    \item $\psi$ -- heading angle (rad)
    \item $v$ -- longitudinal velocity (m/s)
    \item $\delta$ -- virtual bicycle steering angle (rad)
    \item $L$ -- wheelbase (m)
\end{itemize}

\subsection{Derivation of Yaw Rate}

\textbf{Step 1:} Front wheel velocity perpendicular to wheel heading (no-slip):
\begin{equation}
v_f = \frac{v}{\cos(\delta)}
\end{equation}

\textbf{Step 2:} Lateral component of front wheel velocity:
\begin{equation}
v_{f,\text{lateral}} = v_f \sin(\delta) = v \tan(\delta)
\end{equation}

\textbf{Step 3:} Angular velocity about rear axle (instantaneous center of rotation):
\begin{equation}
\dot{\psi} = \frac{v_{f,\text{lateral}}}{L} = \frac{v \tan(\delta)}{L}
\end{equation}

\subsection{Turning Radius}

From the yaw rate equation:
\begin{equation}
R = \frac{L}{\tan(\delta)}
\end{equation}

For small angles ($|\delta| < 0.2$ rad $\approx 11°$):
\begin{equation}
\tan(\delta) \approx \delta \implies R \approx \frac{L}{\delta}
\end{equation}

\subsection{Discrete-Time Integration}

Using Euler forward integration with timestep $\Delta t$:

\begin{align}
\psi_{k+1} &= \psi_k + \Delta t \cdot \frac{v_k}{L} \tan(\delta_k) \\
x_{k+1} &= x_k + \Delta t \cdot v_k \cos(\psi_k) \\
y_{k+1} &= y_k + \Delta t \cdot v_k \sin(\psi_k)
\end{align}

\textbf{Stability condition:} To prevent integration instability:
\begin{equation}
\Delta t < \frac{L}{v_{\max} \tan(\delta_{\max})}
\end{equation}

For $L=2.8$ m, $v_{\max}=60$ m/s, $\delta_{\max}=35° \approx 0.61$ rad:
\begin{equation}
\Delta t < 0.067 \text{ s} = 67 \text{ ms}
\end{equation}

Chosen timestep: $\Delta t = 10$ ms provides 6.7$\times$ safety margin.

\section{Ackermann Steering Geometry}

\subsection{Problem Statement}

Real vehicles have two front wheels that must turn at different angles to maintain pure rolling (no tire scrubbing). All wheels must rotate about a common \textbf{instantaneous center of rotation (ICR)}.

\subsection{Geometric Configuration}

\begin{center}
\begin{tabular}{ll}
\toprule
Parameter & Symbol \\
\midrule
Wheelbase & $L$ \\
Track width (front) & $W$ \\
Virtual steering angle & $\delta$ \\
Turning radius (rear axle to ICR) & $R$ \\
\bottomrule
\end{tabular}
\end{center}

\subsection{Derivation}

\textbf{Step 1:} Turning radius from bicycle model:
\begin{equation}
R = \frac{L}{\tan(\delta)}
\end{equation}

\textbf{Step 2:} Inner and outer wheel radii:
\begin{align}
R_{\text{inner}} &= R - \frac{W}{2} \\
R_{\text{outer}} &= R + \frac{W}{2}
\end{align}

\textbf{Step 3:} Front wheel angles from geometry:
\begin{align}
\delta_{\text{inner}} &= \arctan\left(\frac{L}{R_{\text{inner}}}\right) = \arctan\left(\frac{L}{R - W/2}\right) \\
\delta_{\text{outer}} &= \arctan\left(\frac{L}{R_{\text{outer}}}\right) = \arctan\left(\frac{L}{R + W/2}\right)
\end{align}

\subsection{Turn Direction Convention}

For \textbf{left turn} ($\delta > 0$):
\begin{align}
\delta_{FL} &= \arctan\left(\frac{L}{R - W/2}\right) \quad \text{(inner, larger angle)} \\
\delta_{FR} &= \arctan\left(\frac{L}{R + W/2}\right) \quad \text{(outer, smaller angle)}
\end{align}

For \textbf{right turn} ($\delta < 0$):
\begin{align}
\delta_{FR} &= \arctan\left(\frac{L}{R - W/2}\right) \quad \text{(inner, larger magnitude)} \\
\delta_{FL} &= \arctan\left(\frac{L}{R + W/2}\right) \quad \text{(outer, smaller magnitude)}
\end{align}

\subsection{Small Angle Approximation}

For $|\delta| \ll 1$:
\begin{align}
\delta_{\text{inner}} &\approx \frac{L}{R - W/2} = \delta \left(1 + \frac{W}{2R}\right) \\
\delta_{\text{outer}} &\approx \frac{L}{R + W/2} = \delta \left(1 - \frac{W}{2R}\right)
\end{align}

\section{Longitudinal Dynamics}

\subsection{Force Balance}

Newton's second law in longitudinal direction:
\begin{equation}
m \frac{dv}{dt} = F_{\text{motor}} - F_{\text{brake}} - F_{\text{drag}} - F_{\text{roll}}
\end{equation}

where $m$ is vehicle mass (kg) and $v$ is longitudinal velocity (m/s).

\subsection{Motor Force}

Electric motor torque converted to wheel force:
\begin{equation}
F_{\text{motor}} = \frac{T_{\text{motor}} \cdot N_{\text{gear}} \cdot \eta_{\text{drivetrain}}}{r_{\text{wheel}}}
\end{equation}

where:
\begin{itemize}
    \item $T_{\text{motor}}$ -- motor shaft torque (Nm)
    \item $N_{\text{gear}}$ -- gear ratio (dimensionless)
    \item $\eta_{\text{drivetrain}}$ -- drivetrain efficiency (typically 0.92)
    \item $r_{\text{wheel}}$ -- effective wheel radius (m)
\end{itemize}

\subsection{Power Constraint}

Motors are torque-limited at low speed and power-limited at high speed:
\begin{equation}
P_{\text{motor}} = T_{\text{motor}} \cdot \omega_{\text{motor}}
\end{equation}

Motor angular velocity:
\begin{equation}
\omega_{\text{motor}} = \frac{v}{r_{\text{wheel}}} \cdot N_{\text{gear}}
\end{equation}

Maximum torque at given speed:
\begin{equation}
T_{\text{motor}}^{\max}(v) = \frac{P_{\text{motor}}^{\max}}{\omega_{\text{motor}}(v)} = \frac{P_{\text{motor}}^{\max} \cdot r_{\text{wheel}}}{v \cdot N_{\text{gear}}}
\end{equation}

\subsection{Aerodynamic Drag}

Quadratic drag force:
\begin{equation}
F_{\text{drag}} = \frac{1}{2} \rho C_d A v^2 = C_{\text{drag}} v |v|
\end{equation}

where:
\begin{itemize}
    \item $\rho$ -- air density (\SI{1.225}{kg/m^3} at sea level)
    \item $C_d$ -- drag coefficient (dimensionless, 0.25--1.2 depending on vehicle)
    \item $A$ -- frontal area (m$^2$)
    \item $C_{\text{drag}} = \frac{1}{2}\rho C_d A$ -- combined drag constant
\end{itemize}

\subsection{Rolling Resistance}

Simplified constant friction model:
\begin{equation}
F_{\text{roll}} = C_{\text{roll}} \cdot \text{sign}(v)
\end{equation}

where $C_{\text{roll}}$ is rolling resistance force magnitude (N).

More detailed model:
\begin{equation}
F_{\text{roll}} = \mu_{\text{roll}} \cdot m \cdot g
\end{equation}

where $\mu_{\text{roll}}$ is rolling resistance coefficient (0.01--0.03 for asphalt).

\subsection{Friction Braking}

Brake force proportional to pedal input:
\begin{equation}
F_{\text{brake}} = \frac{p_{\text{brake}}}{100} \cdot \frac{T_{\text{brake}}^{\max}}{r_{\text{wheel}}}
\end{equation}

where $p_{\text{brake}} \in [0, 100]$ is brake pedal percentage.

\subsection{Velocity Update}

Acceleration:
\begin{equation}
a = \frac{F_{\text{motor}} - F_{\text{brake}} - F_{\text{drag}} - F_{\text{roll}}}{m}
\end{equation}

Discrete velocity update:
\begin{equation}
v_{k+1} = v_k + \Delta t \cdot a_k
\end{equation}

\textbf{Zero-crossing logic:} Prevent oscillation near standstill:
\begin{equation}
\text{if } \text{sign}(v_{k+1}) \neq \text{sign}(v_k) \text{ and } |v_{k+1}| < \epsilon_v, \text{ then } v_{k+1} = 0
\end{equation}

where $\epsilon_v = 0.3$ m/s is velocity threshold.

\section{Lateral Dynamics and Friction Limits}

\subsection{Centripetal Acceleration}

During cornering, lateral acceleration is:
\begin{equation}
a_{\text{lat}} = v \cdot \dot{\psi} = \frac{v^2}{L} \tan(\delta)
\end{equation}

\subsection{Friction Circle Constraint}

Maximum lateral acceleration limited by tire-road friction:
\begin{equation}
|a_{\text{lat}}| \leq \mu g
\end{equation}

where:
\begin{itemize}
    \item $\mu$ -- lateral friction coefficient (0.1--0.9 depending on surface)
    \item $g = 9.81$ m/s$^2$ -- gravitational acceleration
\end{itemize}

\subsection{Understeer Implementation}

When friction limit exceeded, scale yaw rate:
\begin{equation}
\dot{\psi}_{\text{actual}} = \dot{\psi}_{\text{commanded}} \cdot \min\left(1, \frac{\mu g}{|a_{\text{lat}}|}\right)
\end{equation}

\section{Battery and Energy Model}

\subsection{State of Charge Dynamics}

Battery state of charge (SOC) evolution:
\begin{equation}
\frac{d(\text{SOC})}{dt} = -\frac{I(t)}{Q_{\text{total}}}
\end{equation}

where:
\begin{itemize}
    \item SOC $\in [0, 1]$ -- normalized state of charge
    \item $I$ -- battery current (A), positive = discharge
    \item $Q_{\text{total}}$ -- total battery capacity (Ah)
\end{itemize}

In terms of power:
\begin{equation}
\frac{d(\text{SOC})}{dt} = -\frac{P(t)}{C_{\text{kWh}} \cdot 1000 \text{ W}}
\end{equation}

where $C_{\text{kWh}}$ is capacity in kWh and $P$ is power in W.

\subsection{Power Flow Directions}

\textbf{Discharge (acceleration):}
\begin{align}
P_{\text{battery}} &= \frac{P_{\text{motor}}}{\eta_{\text{discharge}}} \\
I_{\text{battery}} &= \frac{P_{\text{battery}}}{V_{\text{battery}}} \quad (\text{positive})
\end{align}

\textbf{Charge (regenerative braking):}
\begin{align}
P_{\text{battery}} &= P_{\text{motor}} \cdot \eta_{\text{charge}} \\
I_{\text{battery}} &= \frac{P_{\text{battery}}}{V_{\text{battery}}} \quad (\text{negative})
\end{align}

\subsection{Discrete SOC Update}

Energy flow in timestep $\Delta t$:
\begin{equation}
\Delta E = P \cdot \Delta t \quad \text{(Watt-seconds = Joules)}
\end{equation}

Convert to Watt-hours:
\begin{equation}
\Delta E_{\text{Wh}} = \frac{P \cdot \Delta t}{3600}
\end{equation}

SOC change:
\begin{equation}
\Delta(\text{SOC}) = -\frac{\Delta E_{\text{Wh}}}{C_{\text{kWh}} \cdot 1000}
\end{equation}

Update with saturation:
\begin{equation}
\text{SOC}_{k+1} = \text{clamp}(\text{SOC}_k + \Delta(\text{SOC}), \text{SOC}_{\min}, \text{SOC}_{\max})
\end{equation}

\section{Sensor Noise Models}

\subsection{Gaussian White Noise}

Additive zero-mean Gaussian noise:
\begin{equation}
y = x_{\text{truth}} + \mathcal{N}(0, \sigma^2)
\end{equation}

where $\mathcal{N}(0, \sigma^2)$ is normal distribution with standard deviation $\sigma$.

\textbf{Implementation:} Box-Muller transform:
\begin{align}
U_1, U_2 &\sim \mathcal{U}(0, 1) \quad \text{(uniform random)} \\
Z_1 &= \sqrt{-2 \ln U_1} \cos(2\pi U_2) \\
Z_2 &= \sqrt{-2 \ln U_1} \sin(2\pi U_2)
\end{align}

where $Z_1, Z_2 \sim \mathcal{N}(0, 1)$.

Scaled noise: $\sigma Z_1 \sim \mathcal{N}(0, \sigma^2)$.

\subsection{Random Walk}

Brownian motion (Wiener process):
\begin{equation}
dx = \sigma_{\text{rw}} dW
\end{equation}

where $dW \sim \mathcal{N}(0, dt)$ is Wiener increment.

Discrete implementation:
\begin{equation}
x_{k+1} = x_k + \sigma_{\text{rw}} \sqrt{\Delta t} \cdot \mathcal{N}(0, 1)
\end{equation}

\textbf{Example:} GNSS position drift with $\sigma_{\text{rw}} = 0.1$ m/$\sqrt{\text{s}}$.

\subsection{Gauss-Markov Process}

First-order Markov process (exponentially correlated noise):
\begin{equation}
\frac{db}{dt} = -\frac{1}{\tau} b + w(t)
\end{equation}

where:
\begin{itemize}
    \item $b$ -- bias state
    \item $\tau$ -- correlation time constant (s)
    \item $w(t) \sim \mathcal{N}(0, \sigma_w^2)$ -- white noise driving term
\end{itemize}

\textbf{Discrete solution:}
\begin{equation}
b_{k+1} = e^{-\Delta t / \tau} b_k + \sigma_b \sqrt{1 - e^{-2\Delta t / \tau}} \cdot \mathcal{N}(0, 1)
\end{equation}

where $\sigma_b$ is steady-state bias standard deviation.

\textbf{Steady-state variance:}
\begin{equation}
\sigma_b^2 = \frac{\tau \sigma_w^2}{2}
\end{equation}

\section{IMU Sensor Model}

\subsection{Gyroscope Model}

Measured angular velocity:
\begin{equation}
\omega_{\text{meas}} = \omega_{\text{truth}} + b_{\text{gyro}} + n_{\text{gyro}}
\end{equation}

where:
\begin{itemize}
    \item $\omega_{\text{truth}}$ -- true angular velocity (rad/s)
    \item $b_{\text{gyro}}$ -- time-varying bias (Gauss-Markov, $\tau = 1800$ s)
    \item $n_{\text{gyro}} \sim \mathcal{N}(0, \sigma_{\text{gyro}}^2)$ -- white noise ($\sigma = 0.1°$/s)
\end{itemize}

\textbf{Bias evolution:}
\begin{equation}
b_{\text{gyro},k+1} = \beta b_{\text{gyro},k} + \sigma_b \sqrt{1 - \beta^2} \cdot \mathcal{N}(0, 1)
\end{equation}

where $\beta = \exp(-\Delta t / \tau)$.

\subsection{Accelerometer Model}

Measured acceleration (body frame):
\begin{equation}
\mathbf{a}_{\text{meas}} = \mathbf{a}_{\text{truth}} + \mathbf{b}_{\text{accel}} + \mathbf{n}_{\text{accel}}
\end{equation}

\textbf{Truth values in body frame:}
\begin{align}
a_{x,\text{truth}} &= \frac{dv}{dt} \quad \text{(longitudinal acceleration)} \\
a_{y,\text{truth}} &= v \cdot \dot{\psi} \quad \text{(centripetal acceleration)} \\
a_{z,\text{truth}} &= -g = -9.81 \text{ m/s}^2 \quad \text{(gravity, body Z up)}
\end{align}

\textbf{Bias:} Gauss-Markov with $\tau = 3600$ s (1 hour), $\sigma_b = 0.02$ m/s$^2$.

\textbf{White noise:} $\sigma_{\text{accel}} = 0.05$ m/s$^2$.

\subsection{Gravity Compensation}

Accelerometers measure \textbf{specific force} (excluding gravity). For a stationary vehicle:
\begin{equation}
\mathbf{a}_{\text{meas}} = -\mathbf{g}_{\text{body}} + \text{bias} + \text{noise}
\end{equation}

where $\mathbf{g}_{\text{body}} = [0, 0, -9.81]^\top$ m/s$^2$ in body frame (Z-up).

During motion:
\begin{equation}
\mathbf{a}_{\text{meas}} = \mathbf{a}_{\text{vehicle}} - \mathbf{g}_{\text{body}} + \text{bias} + \text{noise}
\end{equation}

\section{GNSS Sensor Model}

\subsection{Position Measurement}

Latitude and longitude with random walk drift:
\begin{align}
\phi_{\text{meas}} &= \phi_{\text{truth}} + d_{\phi} + \mathcal{N}(0, \sigma_{\phi}^2) \\
\lambda_{\text{meas}} &= \lambda_{\text{truth}} + d_{\lambda} + \mathcal{N}(0, \sigma_{\lambda}^2)
\end{align}

where:
\begin{itemize}
    \item $d_{\phi}, d_{\lambda}$ -- random walk drift ($\sigma_{\text{rw}} = 0.1$ m/$\sqrt{\text{s}}$)
    \item $\sigma_{\phi} = 2.0$ m / $m_{\phi}$ -- white noise (latitude)
    \item $\sigma_{\lambda} = 2.0$ m / $m_{\lambda}$ -- white noise (longitude)
\end{itemize}

\subsection{Velocity Measurement (NED Frame)}

North-East velocity:
\begin{align}
v_{N,\text{meas}} &= v_N + \mathcal{N}(0, 0.1^2) \\
v_{E,\text{meas}} &= v_E + \mathcal{N}(0, 0.1^2)
\end{align}

where $v_N = v \sin(\psi)$ and $v_E = v \cos(\psi)$ from body velocity transformation.

\subsection{Circular Error Probable (CEP)}

The 50\% probability error circle radius:
\begin{equation}
\text{CEP} = 0.59(\sigma_N + \sigma_E) \approx 1.2 \text{ m}
\end{equation}

for $\sigma_N = \sigma_E = 2.0$ m.

\section{Wheel Speed Encoder Model}

\subsection{Quantization Model}

Encoder with $N_{\text{ticks}}$ pulses per revolution:
\begin{equation}
\text{Resolution} = \frac{2\pi}{N_{\text{ticks}}} \text{ rad}
\end{equation}

Quantized measurement:
\begin{equation}
\omega_{\text{quant}} = \left\lfloor \frac{\omega_{\text{truth}}}{\text{Resolution}} \right\rceil \cdot \text{Resolution}
\end{equation}

where $\lfloor \cdot \rceil$ denotes rounding to nearest integer.

\subsection{Total Encoder Noise}

Combined quantization and sensor noise:
\begin{equation}
\omega_{\text{meas}} = \omega_{\text{quant}} + \mathcal{N}(0, \sigma_{\omega}^2)
\end{equation}

where $\sigma_{\omega} = 0.5$ rad/s (white noise).

\subsection{Per-Wheel Variation}

Simulate tire pressure and slip variations:
\begin{equation}
\omega_i = \omega_{\text{meas}} \cdot (1 + u_i)
\end{equation}

where $u_i \sim \mathcal{U}(-0.02, +0.02)$ is uniform random variable (2\% variation).

\section{Radar Sensor Model}

\subsection{Range Measurement}

Simulated stationary target at fixed distance:
\begin{equation}
r_{\text{meas}} = r_{\text{truth}} + \mathcal{N}(0, \sigma_r^2)
\end{equation}

where $r_{\text{truth}} = 50.0$ m and $\sigma_r = 0.2$ m.

\subsection{Doppler Velocity}

Relative radial velocity (closing speed):
\begin{equation}
v_{\text{rel}} = v_{\text{ego}} \cos(\theta)
\end{equation}

where $\theta$ is angle between vehicle heading and line-of-sight to target.

For target dead ahead ($\theta = 0$):
\begin{equation}
v_{\text{meas}} = v_{\text{truth}} + \mathcal{N}(0, \sigma_v^2)
\end{equation}

where $\sigma_v = 0.1$ m/s.

\subsection{Angle Measurement}

Azimuth angle to target:
\begin{equation}
\alpha_{\text{meas}} = \alpha_{\text{truth}} + \mathcal{N}(0, \sigma_{\alpha}^2)
\end{equation}

where $\alpha_{\text{truth}} = 0°$ (dead ahead) and $\sigma_{\alpha} = 0.5°$.

\section{Validation and Verification}

\subsection{Statistical Tests}

\subsubsection{Chi-Squared Goodness of Fit}

Test if noise distribution matches Gaussian:
\begin{equation}
\chi^2 = \sum_{i=1}^{k} \frac{(O_i - E_i)^2}{E_i}
\end{equation}

where $O_i$ are observed counts and $E_i$ are expected counts in histogram bins.

\subsubsection{Allan Variance}

Characterize gyroscope noise:
\begin{equation}
\sigma_A^2(\tau) = \frac{1}{2(N-1)} \sum_{i=1}^{N-1} \left( \bar{\omega}_{i+1}(\tau) - \bar{\omega}_i(\tau) \right)^2
\end{equation}

where $\bar{\omega}_i(\tau)$ is average over interval $\tau$.

\subsection{Coordinate Transform Verification}

\textbf{Round-trip test:} Body $\to$ NED $\to$ Body should preserve magnitude:
\begin{equation}
v = \sqrt{v_N^2 + v_E^2} \quad \text{(Pythagorean theorem)}
\end{equation}

\textbf{Gravity check:} Stationary vehicle should measure:
\begin{equation}
\mathbf{a}_{\text{meas}} = [0, 0, -9.81]^\top + \text{small noise/bias}
\end{equation}

\section{Conclusion}

This mathematical framework provides the foundation for:
\begin{itemize}
    \item High-fidelity vehicle dynamics simulation
    \item Realistic sensor noise modeling for ML/AI training
    \item Coordinate transformations between body, NED, and WGS84 frames
    \item Extended Kalman Filter (EKF) development and tuning
    \item Sensor fusion algorithm validation
\end{itemize}

All models are validated against industry standards and real sensor specifications.

\begin{thebibliography}{99}

\bibitem{rajamani2012}
Rajamani, R. (2012). \textit{Vehicle Dynamics and Control}. Springer.

\bibitem{gillespie1992}
Gillespie, T. D. (1992). \textit{Fundamentals of Vehicle Dynamics}. SAE International.

\bibitem{groves2013}
Groves, P. D. (2013). \textit{Principles of GNSS, Inertial, and Multisensor Integrated Navigation Systems}. Artech House.

\bibitem{ieee1554}
IEEE Standard 1554-2005. \textit{Recommended Practice for Inertial Sensor Test Equipment, Instrumentation, Data Acquisition, and Analysis}.

\bibitem{polack2017}
Polack, P., et al. (2017). "The Kinematic Bicycle Model: A Consistent Model for Planning Feasible Trajectories for Autonomous Vehicles?" \textit{IEEE Intelligent Vehicles Symposium}.

\bibitem{woodman2007}
Woodman, O. J. (2007). "An introduction to inertial navigation." University of Cambridge, Computer Laboratory, Tech. Rep. UCAM-CL-TR-696.

\end{thebibliography}

\end{document}